\frontsection{Приложение Б. Шаблон отчёта}
\label{app:report-template}
Шаблон задаёт структуру содержания. Требования к оформлению титульного листа, объёму и сроку сдачи сообщает преподаватель. Заполните поля своими сведениями; таблицы и графики включайте по смыслу исследования.

\subsection*{Сведения о работе}
\noindent Дисциплина: \hrulefill\par
\noindent Лабораторная работа №\,\rule{8mm}{0.4pt}. Название: \hrulefill\par
\noindent Выполнил(а): \hrulefill\quad Группа: \rule{25mm}{0.4pt}\par
\noindent ИТМО ID: \rule{28mm}{0.4pt}\quad Проверил(а): \hrulefill\par
\noindent Дата: \rule{28mm}{0.4pt}\quad Версия материалов: \hrulefill

\subsection*{1. Цель и данные}
Сформулируйте цель своими словами. Укажите источник и версию файлов, объект наблюдения, признаки и целевой ответ. Приведите размер набора и проверки его состава. Объясните, какие данные доступны при применении модели.

\subsection*{2. Условия эксперимента}
Опишите разбиение и способ исключения утечки. Укажите среду, seed, модель, основную метрику выбора и бюджет. Для каждого опыта сформулируйте ожидаемый эффект изменения параметра до запуска. Дополнительные опыты явно отделяйте от исходного плана.

\subsection*{3. Результаты и разбор ошибок}
Используйте таблицу следующего вида; численные ячейки заполняются после выполнения опытов.
\begin{center}\small
\begin{tabularx}{\textwidth}{@{}p{0.13\textwidth}Y Y Y@{}}
\toprule
Опыт & Что изменено & Метрика на указанной части & Наблюдение \\
\midrule
Исходный & Конфигурация модели & Значение и единицы & Что видно по результату \\
Сравнение & Один фактор либо явно заданная комбинация & Та же метрика и часть & Изменение относительно исходного \\
\bottomrule
\end{tabularx}
\end{center}
Сошлитесь на диагностические графики и разберите конкретные ошибки. При неудачном запуске укажите причину, а не подставляйте нулевое значение метрики.

\subsection*{4. Выводы и приложение к отчёту}
Ответьте, достигнута ли цель, что подтверждено результатами и чем ограничен вывод. Приложите ноутбук, конфигурации, таблицы опытов и необходимые файлы модели. Перед сдачей проверьте выполнение ноутбука сверху вниз и соответствие его результатов таблицам отчёта.
